\documentclass[11pt]{article}
\usepackage[margin=1in]{geometry}
\usepackage{amsmath,amssymb,amsthm}
\usepackage{hyperref}
\usepackage{xurl}
\let\oldus\_
\renewcommand{\_}{\oldus\allowbreak}

\newtheorem{theorem}{Theorem}
\newtheorem{lemma}[theorem]{Lemma}
\theoremstyle{definition}
\newtheorem{definition}[theorem]{Definition}
\theoremstyle{remark}
\newtheorem{remark}[theorem]{Remark}

\title{A disproof of Conjecture 00000002221\\[2pt]
\large $\mathbb F_2^{\,3}$ has Krull dimension $0$ but a Stone space with $3>2^{2^0}$ points}
\date{}

\begin{document}
\sloppy
\maketitle

\section{The conjecture and its readings}

Conjecture 00000002221 reads: \emph{Definition: The constructible topology on $\operatorname{Spec} R$.
Conjecture: The cardinality of the Stone space of the Booleanization (complemented-closure algebra) of the
spectrum is at most $2^{2^d}$ with $d=\dim$, explicitly (Booleanization cardinality).}

Here $R$ is a commutative ring, $d=\dim R$ is its Krull dimension (equal to the Krull dimension of the
poset $\operatorname{Spec} R$), and the Stone space of a Boolean algebra $B$ is the set of Boolean
homomorphisms $B\to\mathbf 2=\{0,1\}$ (equivalently, the ultrafilters of $B$).

The phrase ``Booleanization (complemented-closure algebra) of the spectrum'' is not a standard term. We
refute the following readings of it.
\begin{enumerate}
\item[(R1)] $B$ is the Boolean algebra of \emph{constructible} subsets of $\operatorname{Spec} R$: the
Boolean subalgebra of the power set generated by the retrocompact open sets (Mathlib's
\texttt{Topology.IsConstructible}). This is the reading suggested by the ``Definition'' line, because the
Stone space of this algebra is $\operatorname{Spec} R$ with the constructible topology.
\item[(R2)] $B$ is the Boolean algebra of \emph{clopen} subsets of $\operatorname{Spec} R$, i.e.\ the
complemented elements of the lattice of open sets.
\item[(R3)] The ``Stone space'' is taken to be $\operatorname{Spec} R$ itself with the constructible
topology.
\item[(R4)] $B$ is any Boolean algebra with a homomorphism of bounded lattices $\varphi:B\to
\mathcal P(\operatorname{Spec} R)$ whose image contains every singleton. This includes the full power set.
\end{enumerate}
For the ring used below, every subset of $\operatorname{Spec} R$ is both clopen and constructible
(Lemma~\ref{lem:subsets}). So the algebras in (R1) and (R2) are the full power set, and both are instances
of (R4). Readings in which $B$ is not a Boolean algebra of subsets of $\operatorname{Spec}R$ are not
treated.

\section{The counterexample}

Let $K$ be a field, $N\ge1$, and $R=K^N=\prod_{i<N}K$. For $i<N$ let $\mathfrak p_i=\ker(\mathrm{ev}_i)$,
where $\mathrm{ev}_i:R\to K$ is the $i$-th projection.

\begin{lemma}
$\dim R=0$.
\end{lemma}
\begin{proof}
$R$ is a finite-dimensional $K$-algebra, hence Artinian, so every prime ideal is maximal. Since $R\ne0$,
$\operatorname{Spec}R\neq\emptyset$, and therefore $\dim R=0$.
\end{proof}

\begin{lemma}\label{lem:subsets}
The map $\coprod_{i<N}\operatorname{Spec}K\to\operatorname{Spec}R$, $(i,\mathfrak q)\mapsto
\mathrm{ev}_i^{-1}(\mathfrak q)$, is a homeomorphism. Consequently:
\begin{itemize}
\item $\operatorname{Spec}R$ is a finite discrete space;
\item the primes $\mathfrak p_0,\dots,\mathfrak p_{N-1}$ are pairwise distinct;
\item every subset of $\operatorname{Spec}R$ is clopen and constructible.
\end{itemize}
\end{lemma}
\begin{proof}
The homeomorphism is Mathlib's \texttt{PrimeSpectrum.sigmaToPiHomeo}, valid for finite products.
$\operatorname{Spec}K$ is a single point, so the coproduct is discrete and finite, and the
$\mathfrak p_i$ are its distinct images. In a discrete space every set is open and closed. In a finite
space every set is compact, so every open set is retrocompact. Hence every set is a retrocompact open, and
in particular constructible.
\end{proof}

\begin{lemma}\label{lem:stone}
Let $X$ be a set, $B$ a Boolean algebra, and $\varphi:B\to\mathcal P(X)$ a homomorphism of bounded
lattices whose image contains every singleton. Then $x\mapsto\bigl(b\mapsto[x\in\varphi(b)]\bigr)$ is an
injective map from $X$ into the Stone space $\operatorname{Hom}(B,\mathbf 2)$. Hence
$|\operatorname{Hom}(B,\mathbf 2)|\ge|X|$.
\end{lemma}
\begin{proof}
Evaluation at a point preserves $\cup,\cap,\emptyset,X$, so each map is a Boolean homomorphism. If
$x\ne y$, choose $b$ with $\varphi(b)=\{x\}$; the two homomorphisms differ at $b$.
\end{proof}

\begin{theorem}
For every field $K$ and every $N\ge1$, $\dim K^N=0$. In each of the readings (R1)--(R4), the Stone space
of $\operatorname{Spec}K^N$ has at least $N$ points. In particular, for $R=\mathbb F_2^{\,3}$ we have
$d=0$ and the bound would be $2^{2^0}=2$, but the Stone space has at least $3$ points. So the conjecture is
false in readings (R1), (R2) and (R3), and in (R4) for every admissible $B$.
\end{theorem}

\begin{remark}
In fact $\operatorname{Spec}K^N$ has exactly $N$ points, and the Stone space of the full power set of an
$N$-point set has exactly $N$ points. Only the lower bound is needed and formalized. Even if ``cardinality''
were applied to the Boolean algebra itself ($2^3=8$ elements), the bound $2$ would fail.
\end{remark}

\section{Relation to the earlier submission}
PR \#219 by orionsheep was closed unmerged. The reason was that its Lean ``proves only $2^{(2^0)} = 2$,
$3 > 2$, $8 > 2$; the ring $k^3$, Spec, Booleanization, Stone space, and Krull dimension appear nowhere.''
Its prose also counted the Stone space as having $2^3=8$ points. That is the cardinality of the Boolean
algebra $\mathcal P(\operatorname{Spec}k^3)$, not of its Stone space, which has $3$ points. The
conclusion survives, since $3>2$ as well.

The present Lean file works with the actual objects:
\begin{itemize}
\item the ring \texttt{Fin N $\to$ K} and its Krull dimension \texttt{ringKrullDim};
\item \texttt{PrimeSpectrum} with its Zariski topology;
\item the constructible Boolean algebra, as Mathlib's \texttt{BooleanSubalgebra.closure} of the
retrocompact opens;
\item the clopen algebra \texttt{TopologicalSpace.Clopens};
\item Mathlib's \texttt{WithConstructibleTopology};
\item the Stone space \texttt{BoundedLatticeHom B Bool}.
\end{itemize}
The main theorems quantify over all commutative rings $R$ and refute the bound for each reading.

\section{Lean formalization}
The file is \texttt{lean/Conjecture2221/Basic.lean} (namespace \texttt{C2221}). It imports only Mathlib.
\begin{itemize}
\item Definitions:
  \begin{itemize}
  \item \texttt{StoneSpace B := BoundedLatticeHom B Bool};
  \item \texttt{constructibleAlgebra R := BooleanSubalgebra.closure \{U | IsOpen U $\wedge$
  IsRetrocompact U\}} (membership is \texttt{IsConstructible}, see
  \texttt{mem\_constructibleAlgebra\_iff});
  \item \texttt{clopensHom}, the inclusion of clopens into sets;
  \item \texttt{evalAt}, the evaluation ultrafilters $b\mapsto[x\in\varphi(b)]$ (principal for the finite power-set algebras of R1 and R2; for a general $B$ in R4 they need not be principal, and only evaluation and separation are used);
  \item \texttt{bound d := 2\textasciicircum 2\textasciicircum d}.
  \end{itemize}
\item \texttt{evalAt\_injective} and \texttt{mk\_le\_stone}: Lemma~\ref{lem:stone}.
\item \texttt{ringKrullDim\_pi : ringKrullDim (Fin N $\to$ K) = 0} for $N\ne0$.
\item Facts about $\operatorname{Spec}$:
  \begin{itemize}
  \item instances \texttt{discrete\_spec} and \texttt{finite\_spec};
  \item \texttt{pt\_injective} and \texttt{card\_spec}: $N\le\#\operatorname{Spec}$;
  \item \texttt{every\_subset}: every subset is \texttt{IsClopen} and \texttt{IsConstructible}.
  \end{itemize}
\item Lower bounds $N\le\#$ for each reading:
  \begin{itemize}
  \item \texttt{stone\_card\_of\_hom} (R4);
  \item \texttt{stone\_constructible\_card} (R1);
  \item \texttt{stone\_clopens\_card} (R2);
  \item \texttt{constructibleTopology\_card} (R3).
  \end{itemize}
\item Main theorems \texttt{conjecture\_false\_constructible}, \texttt{conjecture\_false\_clopens} and
\texttt{conjecture\_false\_constructibleTopology}. Each has the form
\[
\neg\,\forall\,(R:\texttt{Type})\,[\texttt{CommRing }R]\,(d:\mathbb N),\
\texttt{ringKrullDim }R=d\to \#(\text{Stone space})\le 2^{2^d},
\]
and is witnessed by $R=$ \texttt{Fin 3 $\to$ ZMod 2}.
\item Axiom audit: \texttt{\#print axioms} for the three main theorems, \texttt{stone\_card\_of\_hom} and
\texttt{every\_subset} reports \texttt{[propext, Classical.choice, Quot.sound]}. There is no
\texttt{sorry} and no \texttt{native\_decide} (the ordinary \texttt{decide} term is used once, inside \texttt{evalAt}, to decide membership $x\in\varphi(b)$).
\end{itemize}

\end{document}
